% Part: applied-modal-logics
% Chapter: temporal-logic
% Section: possible-histories

\documentclass[../../../include/open-logic-section]{subfiles}

\begin{document}

\olfileid{aml}{tl}{poss}

\olsection{تاریخ‌های ممکن}

مدل‌های رابطه‌ایِ منطق زمانی که تاکنون به کار برده‌ایم بسیار انعطاف‌پذیر
هستند، زیرا لازم نیست هیچ محدودیتی بر رابطهٔ دسترسی‌پذیری بگذاریم. این
یعنی مدل‌های زمانی می‌توانند هم در گذشته و هم در آینده شاخه‌دار باشند؛
اما شاید بخواهیم تصوری «موجهاتی‌تر» از شاخه‌شدن را در نظر بگیریم که در
آن دنباله‌های رویدادها را تاریخ‌های ممکن می‌دانیم. این امر لزوماً
نیازمند تغییر زبان نیست، هرچند ممکن است عملگرهای موجهاتی «معمول»مان،
$\Box$ و~$\Diamond$، را نیز بیفزاییم؛ و حتی افزودن روابط دسترسی‌پذیری
معرفتی را برای بازنمایی تغییرات دانش عامل‌ها در طول زمان در نظر بگیریم.

\begin{defn}
  یک \emph{مدل تاریخ‌های ممکن} برای زبان زمانی سه‌تاییِ
  $\mModel{M} = \tuple{T, C, V}$ است، که در آن:
  \begin{enumerate}
    \item $T$ مجموعه‌ای ناتهی است که به‌عنوان حالت‌های زمانی تفسیر
      می‌شود.
    \item $C$ مجموعه‌ای از مسیرهای محاسباتی، یا \emph{تاریخ‌های ممکن}
      یک دستگاه است. به بیان دیگر، $C$ مجموعه‌ای از دنباله‌های~$\sigma$
      از حالت‌های~$s_1$، $s_2$،~$s_3$، \dots{} است که هر~$s_i \in T$.
    \item $V$ تابعی است که به هر !!{propositional variable}~$p$
      مجموعه‌ای~$V(p)$ از نقاط زمانی تخصیص می‌دهد.
  \end{enumerate}
  برای ساده‌ترشدن امور، عموماً فرض می‌کنیم هرگاه تاریخی عضو~$C$ باشد،
  همهٔ پسوندهای آن نیز عضو~$C$ باشند. برای نمونه، اگر~$s_1$، $s_2$،
  $s_3$ دنباله‌ای در~$C$ باشد، آنگاه~$s_2$، $s_3$ و نیز~$s_3$ هم
  در~$C$ هستند. همچنین وقتی دو حالت~$s_i$ و~$s_j$ در دنبالهٔ~$\sigma$
  ظاهر شوند، هرگاه~$i < j$ می‌گوییم~$s_i \prec_\sigma s_j$. وقتی
  $t \in V(p)$، می‌گوییم~$p$ \emph{در}~$t$ \emph{صادق است}.
\end{defn}

تغییر مربوط این است که وقتی صدق !!a{formula} را در نقطهٔ زمانی~$t$ در
مدل~$\mModel M$ ارزیابی می‌کنیم، این کار را نسبت به تاریخ~$\sigma$ای
انجام می‌دهیم که~$t$ در آن به‌صورت یک حالت ظاهر می‌شود. بااین‌حال،
لازم نیست هیچ‌یک از معناشناسی‌های !!{propositional variable}ها یا
رابط‌های تابع‌صدقی را تغییر دهیم. همهٔ آن‌ها دقیقاً همانند
\olref[sem]{defn:tmodels} هستند، زیرا هیچ‌یک به~$\sigma$ ارجاع نمی‌دهند.
اما اکنون عملگر آیندهٔ~$\Ftemp$ را بازتعریف می‌کنیم و عملگر~$\Diamond$
را نسبت به این تاریخ‌ها می‌افزاییم.

\begin{defn}\ollabel{defn:phmodels}
  \emph{صدق !!a{formula}~$!A$ در~$t, \sigma$} در~$\mModel M$، با نماد
  $\mSat{M}{!A}[t, \sigma]$:
  \begin{enumerate}
  \item\ollabel{defn:sub:phmodels-f} \indcase{!A}{\Ftemp
    !B}{$\mSat{M}{\indfrm}[t, \sigma]$ اگر و تنها اگر برای یک
    $t' \in T$ چنان‌که~$t \prec_\sigma t'$،
    $\mSat{M}{!B}[t', \sigma]$.}
  \item\ollabel{defn:sub:phmodels-diamond} \indcase{!A}{\Diamond
    !B}{$\mSat{M}{\indfrm}[t, \sigma]$ اگر و تنها اگر برای یک
    $\sigma' \in C$ که~$t$ در آن ظاهر می‌شود،
    $\mSat{M}{!B}[t, \sigma']$.}
  \end{enumerate} 
\end{defn}

می‌توان دیگر عملگرهای زمانی و موجهاتی را نیز به‌همین‌سان تعریف کرد.
بااین‌حال، اکنون می‌توانیم ادعاهایی را بازنمایی کنیم که زمان دستوری و
وجه را ترکیب می‌کنند. برای نمونه، می‌توانیم «$p$ رخ نخواهد داد، اما
ممکن بود رخ دهد» را با فرمول~$\lnot \Ftemp p \land \Diamond \Ftemp p$
نمادسازی کنیم. این فرمول در نقطه و تاریخی برقرار است که در آن~$p$ در
هیچ حالت پسینی صادق نمی‌شود، اما تاریخی بدیل وجود دارد که در آن~$p$
صادق خواهد شد.

\end{document}
